% 第八章 总结与展望

\chapter{总结与展望}


\section{全文总结}

本文围绕 FreeFlow 原型系统，完成了一个面向链上音乐发布、授权访问与播放消费的 Web2.5 音乐平台原型的设计与实现；设计了一套冷启动场景下的曲目检索方法，并完成在智能合约功能性、检索方法的有效性在模拟数据集下的测试。



\section{不足与改进方向}

当前系统仍存在一些不足之处。当前FreeFlow中音频资源通过Pinata 网关访问，截获CID后仍可能绕过系统访问控制，后续可引入音频加密、授权网关等服务解决； 目前检索排序模块目前主要基于元数据、链上身份字段和聚合热度，缺少真实用户点击、播放和收藏反馈，且评估基于MusicBrainz 样本库，无法反映大规模生产场景下用户行为；智能合约基于Hardhat框架离线运行，仍需考虑真实链上gas成本优化。
此外，鉴于复杂度和时间可行性权衡，完整社区治理能力尚未实现。评论、检索索引等高频业务数据仍然通过中心化服务器管理,若要通过治理代币解决内容审核和自建IPFS存储，仍需要后续继续扩展。
